\documentclass[10pt,a4paper]{amsart}
\usepackage[margin=19mm]{geometry}
\usepackage{amsmath,amssymb,mathtools}
\usepackage[scr=rsfs,cal=euler]{mathalfa}
\usepackage{xfrac}
\usepackage[unicode,hidelinks,bookmarks=false]{hyperref}
\newcommand{\ie}{i.e.\ }
\newcommand{\cf}{cf.\ }
\newcommand{\resp}{respectively\ }
\newcommand{\Subsection}{\subsection}
\providecommand{\nobreakdash}{\nobreak\hbox{-}}
\setlength{\emergencystretch}{2em}
\multlinegap0pt
\usepackage{amssymb}
\usepackage{enumerate}
\usepackage{tikz}
\usepackage{bm}
\usepackage{float}
\newtheorem{lemma}{Lemma}
\newtheorem{theorem}{Theorem}
\title{The second-class particle in the half-line open TASEP}
\author{Kailun Chen}
\date{Source: arXiv:2409.00554v4 (CC BY 4.0)}
\begin{document}
\maketitle

\noindent\emph{Source and licence.} The second-class particle in the half-line open TASEP. Version arXiv:2409.00554v4 (CC BY 4.0), distributed by arXiv under the Creative Commons Attribution 4.0 International licence, http://creativecommons.org/licenses/by/4.0/, as recorded in the arXiv licence metadata for this exact version and shown on the abstract page https://arxiv.org/abs/2409.00554v4. Full licence notice: \texttt{licenses/arxiv-2409.00554-CC-BY-4.0.txt}.\par\smallskip

\noindent\textit{Editorial scope.} Counted unit: the article's theorem thm:color-position (source0.tex lines 920--923). The article's complete original proof of it is delivered with it (source0.tex lines 925--977, one \texttt{proof} environment of 53 lines). No mathematics is written, repaired or re-derived: the statement and the proof are the article's own lines, in the article's own notation, and no retained line is substituted or re-flowed.  Retained uncounted as the source-local context the proof needs: lines 805--821; lines 862--868.  Nothing was omitted from any retained span, so the package carries no omission record.  No retained line contains a citation, so the wrapper carries no bibliography and no bibliography entry.  No retained line contains a figure environment, an image-inclusion command or drawing code, so no image is delivered and the package declares no asset dependency.  Editorial formatting only, in addition: an AMS \texttt{amsart} preamble that restates the article's own declarations and macros used by the retained lines, namely the environments \texttt{lemma}, \texttt{theorem} on separate counters, together with the standard packages those lines need.  The retained environments print in this document as Lemma 1, Theorem 1 in order of appearance, whereas the article prints the same items with the numbering of its own full document.  The complete native source of the article as distributed in the arXiv e-print for this exact version is bundled unchanged as \texttt{sources/164570/source0.tex}.
\begin{lemma}
\label{lem:clock-sequence-hecke-product}
Fix a realization of the clocks up to time $t$ for which only the clocks indexed by
$a_1,a_2,\ldots,a_r$ ring. 
Then the colored half-space TASEP configuration at time $t$, started from $\eta_0\in\Omega_n$, is
\[
U_{a_r}\cdots U_{a_2}U_{a_1}\eta_0.
\]
Equivalently, under the identification of configurations with basis elements of the type $B/C$ Hecke algebra, this configuration is represented by
\[
T_{s_{a_r}}\cdots T_{s_{a_2}}T_{s_{a_1}}T_{\eta_0}.
\]
In particular, for the packed initial condition $\eta_0=id$, the random walk $W(t)$ satisfies
\[
W(t)=T_{s_{a_r}}\cdots T_{s_{a_2}}T_{s_{a_1}}.
\]
\end{lemma}

\begin{lemma}
\label{lem:d-equality}
With $\overset{\text{d}}{=}$ denoting equality in distribution, we have 
\begin{align}
\iota(W(t)) \overset{\text{d}}{=} W(t).
\end{align}
\end{lemma}

\begin{theorem}[color--position symmetry]
\label{thm:color-position}
The random configurations $\eta_t^{\mathrm{gen};1}$ and $\iota(\eta_t^{\mathrm{gen};2})$ have the same distribution. Equivalently, $\iota(\eta_t^{\mathrm{gen};1})$ and $\eta_t^{\mathrm{gen};2}$ have the same distribution.
\end{theorem}

\begin{proof}
We prove the first distributional identity; the second follows by applying the same argument after taking inverses.

Condition on a realization of the clocks up to time $t$, and write the corresponding Hecke word as
\[
W(t)=T_{s_{a_r}}\cdots T_{s_{a_1}}.
\]
By Lemma \ref{lem:clock-sequence-hecke-product} and by the definitions of the two generated processes,
\[
\eta_t^{\mathrm{gen};1}
=
T_{s_{a_r}}\cdots T_{s_{a_1}}
T_{s_{k_n}}\cdots T_{s_{k_1}},
\]
whereas
\[
\eta_t^{\mathrm{gen};2}
=
T_{s_{k_1}}\cdots T_{s_{k_n}}
T_{s_{a_r}}\cdots T_{s_{a_1}}.
\]
Applying the anti-involution $\iota$ and using $\iota(T_{s_j})=T_{s_j}$ for every Coxeter generator gives
\[
\iota(\eta_t^{\mathrm{gen};2})
=
\iota\!\left(
T_{s_{k_1}}\cdots T_{s_{k_n}}
T_{s_{a_r}}\cdots T_{s_{a_1}}
\right)
=
\iota(W(t))\,T_{s_{k_n}}\cdots T_{s_{k_1}}.
\]
By Lemma \ref{lem:d-equality}, $\iota(W(t))\overset{d}{=}W(t)$. 
Therefore
\[
\iota(\eta_t^{\mathrm{gen};2})
\overset{d}{=}
W(t)T_{s_{k_n}}\cdots T_{s_{k_1}}
=
\eta_t^{\mathrm{gen};1}.
\]
This proves
\[
\eta_t^{\mathrm{gen};1}\overset{d}{=}\iota(\eta_t^{\mathrm{gen};2}).
\]

Finally, since $\iota$ is an involution, applying $\iota$ to both sides gives
\[
\iota(\eta_t^{\mathrm{gen};1})
\overset{d}{=}
\eta_t^{\mathrm{gen};2}.
\]
\end{proof}

\end{document}
